\documentclass[a4paper,12pt]{report}
\newcommand{\path}{../}
\usepackage[fiche,niveau={1MA1},nfiche={22}, annee={Emilie-Gourd, 2025--2026},
auteur={ns},theme={Géométrie -- Preuves du cercle de Thalès}]{\path packages/bandeaux}
\usepackage{\path packages/boites}
\usepackage{\path packages/courslatex}
\usepackage{\path packages/fontes}
\usepackage{\path packages/mathenv}
\allowdisplaybreaks
\begin{document}
\begin{thm}[Théorème du cercle de Thalès]
	Un triangle inscrit dans un cercle et dont un côté est un diamètre est un triangle rectangle. 
\end{thm}
\begin{proof}
	\begin{minipage}[t]{0.5\textwidth}{
	\vspace{0pt}
	Le triangle $AOC$ est isocèle en $O$, car $AO$ et $CO$ sont des rayons. Ainsi, $\alpha=\widehat{ACO}=\widehat{CAO}$. 

Le triangle $BOC$ est isocèle en $O$, car $BO$ et $CO$ sont des rayons. Ainsi, $\beta=\widehat{BCO}=\widehat{CBO}$.

De plus, par la somme des angles du triangle $ACB$, on a

\[\alpha+\alpha+\beta+\beta=180^\circ\]
 
d'où 
\[\alpha+\beta=90^\circ.\]
Ainsi, $\widehat{ACO}+\widehat{BCO}=90^\circ$ et le triangle $ABC$ est rectangle en $C$. 
	}
	\end{minipage}
	\begin{minipage}[t]{0.4\textwidth}{
	\vspace{0pt}
\begin{center}
	\begin{tikzpicture}[scale=1.5]
\tkzSetUpPoint[shape=cross]
  \tkzInit[xmin=-2,xmax=2,ymin=-2,ymax=2]
  \tkzClip
  \tkzDefPoints{-1/0/A, 1/0/B}           % Define points A and B
  \tkzDefMidPoint(A,B) \tkzGetPoint{O}   % Midpoint O is the center of the circle
  \tkzDefPointBy[rotation=center O angle 60](B) \tkzGetPoint{C}  % C on the circle, forming a right triangle
  \tkzDrawCircle(O,A)                    % Draw the circle through A with center O
  \tkzDrawPolygon(A,C,B)                 % Draw triangle ABC
  \tkzDrawPoints(A,B,C,O)
  \tkzDrawSegment[dashed](O,C)
  \tkzLabelPoints(O)
  \tkzLabelPoints[above](C)
  \tkzLabelPoints[right](B)
  \tkzLabelPoints[left](A)
\end{tikzpicture}
\end{center}
	}
	\end{minipage}

\end{proof}
\begin{thm}[Réciproque du théorème du cercle de Thalès]
Si un triangle $ABC$ inscrit dans un cercle est rectangle en $B$ alors $AC$ est un diamètre du cercle.  
\end{thm}
\begin{proof}
\begin{minipage}[t]{0.5\textwidth}{
\vspace{0pt}
Soient $O$ le milieu du côté $AC$ et $A'$ le milieu du côté $CB$. On trace la droite $OA'$ passant par $O$ et $A'$. Par la réciproque du théorème de Thalès, $d_{OA'}\| d_{AB}$. Ainsi, $d_{OA'}\perp d_{CB}$ et $d_{OA'}$ est la médiatrice du côté $BC$.
\medskip

Soit $C'$ le milieu du côté $AB$. On trace la droite $OC'$ passant par $O$ et $C'$. Par la réciproque du théorème de Thalès, $d_{OC'}\| d_{CB}$. Ainsi, $d_{OC'}\perp d_{AB}$ et $d_{OC'}$ est la médiatrice du côté $AB$.
\medskip

Les deux médiatrices $d_{OA'}$ et $d_{OC'}$ s'intersectent au point $O$ qui est le centre du cercle circonscrit dont le côté $AC$ est le diamètre (car c'est l'hypoténuse). 
}
\end{minipage}
\hfill
\begin{minipage}[t]{0.4\textwidth}{
\vspace{0pt}
\begin{tikzpicture}[scale=0.75]

\tkzSetUpPoint[shape=cross]
  % Define the vertices of triangle ABC
  \tkzDefPoint(0,0){A}
  \tkzDefPoint(6,0){B}
  \tkzDefPoint(6,8){C}
  
  % Draw the triangle
  \tkzDrawPolygon(A,B,C)
  \tkzDrawPoints(A,B,C)
  \tkzLabelPoints[below left](A)
  \tkzLabelPoints[below right](B)
  \tkzLabelPoints[above](C)
  
  % Compute the midpoints:
  % O is the midpoint of AC
  \tkzDefMidPoint(A,C) \tkzGetPoint{O}
  % A' is the midpoint of BC
  \tkzDefMidPoint(B,C) \tkzGetPoint{A'}
  % C' is the midpoint of AB
  \tkzDefMidPoint(A,B) \tkzGetPoint{C'}
  
  \tkzDrawPoints(O,A',C')
  \tkzLabelPoints[above left](O)
  \tkzLabelPoints[above right](A')
  \tkzLabelPoints[below left](C')
  
  % Draw the perpendicular bisectors (mediatrices)
  \tkzDrawLine[dashed, add=0.5 and 0.5](O,A')
  \tkzDrawLine[dashed, add=0.5 and 0.5](O,C')
  
  % Mark the right angles (showing perpendicularity)
  \tkzMarkRightAngle[scale=0.5](B,A',O)
  \tkzMarkRightAngle[scale=0.5](A,C',O)
  \tkzMarkRightAngle(A,B,C)
  
  % Draw the circle with diameter AC (circumcircle)
  \tkzDrawCircle[dashed](O,A)
\end{tikzpicture}

}
\end{minipage}

\end{proof}

\end{document}
